\section{Induction} $ $\\

Is a tomato a fruit? Perhaps.
\newline
Who decides what isn't a fruit and what is? Botanists may say that tomatoes are a fruit because they from a flower and contain seeds. Botanists don't care how tomatoes taste in a fruit salad. 
\newline
When we are defining combinatorial structures and talking about them, we care about aspects of these objects that perhaps others using them are not particularly interested in. Botanists define fruit in a way that make it very clear what is and isn't a fruit (according to them anyway).
\newline
The following conversation regarding the construction of the positive natural numbers is a bit like botanists discussing "what is a fruit and how do we define it".
\newline

Virtually all of the infinite sets that we have constructed so far have been based off the existence of a well understood infinite set, the natural numbers. We have not yet formally defined or constructed them. Is there only one set of Natural numbers? Are there many? Is it up to interpretations? Below are a handful of observations that led to the construction we all use today.
\newline

We have a pretty good idea of what the positive natural numbers are $1,2,\cdots $, but what the hell does the $\cdots$ mean? Is infinity a natural number then? $\infty\in \mathbb{N}$? Are there an infinite number of natural numbers?
\newline
We do have an idea of a next natural number, we define it as the successor.
\begin{defn}{successor}
If $n$ is a positive integer, we call $S(n)$ the successor of $n$ and we say $n<S(n)$. We write $S(n)=n+1$.
\end{defn}
We also have an idea of what it means for one natural number to be bigger than another. 
\begin{defn}{$<$} $ $ \\
We define the binary operation $<:\mathfrak{U}\times \mathfrak{U}\to \{\text{True},\text{False}\}$ where $ a < S(a)$ and $<$ satisfies the transitive property $a<b\land b<c \to a <c$. If $a < b$ we say $a$ is less than $b$.
\end{defn}
Now we can try to define the natural numbers:
\begin{exmp}[Constructing infinite sets]
Consider a set $N$ with the following properties:
\begin{enumerate}
\item It contains $1$
\item If $n$ is in the set, then the successor of $n$ is a natural number.
\item All positive integers are the successor of a unique positive integer or are $1$.
\end{enumerate}
We can imagine the members of these sets as nodes, or dots. We can draw an arrow from a dot $a$ to a dot $b$ if the successor of $a$ is $b$.
\newline
Do we all have the same drawing? Is a set defined this way unique? or can our drawings look different? Can you prove to me that there is no Natural number that ends in a $.7$? It doesn't seem like there should be one?
\newline
\vspace{2cm}
\end{exmp}


\newpage
As it turns out, this is not enough. We can have all sorts of weird looking collections of dots that don't break any of these rules. We need one additional rule:

\begin{defn}[The well ordering principle]
Every nonempty subset of the set of positive
integers has a least element.
\end{defn}
\begin{exmp} Prove that no natural number ends in a $.7$.
\vspace{3cm}
\end{exmp}


As it turns out, the well ordering principle is equivalent to the axiom of mathematical induction.

\begin{defn}[Axiom of Mathematical Induction]
If $S$ is a set of positive integers such that $1 \in S$ and
for all positive integers $n$ if $n \in S$, then $n + 1 \in S$, then $S$ is the set of positive integers.
\end{defn}
We can use mathematical induction to prove the above example as well.
\begin{exmp} Prove that no natural number ends in a $.7$.
\vspace{3cm}
\end{exmp}



\begin{defn}[Induction] $ $\\
Inference Rule:
\newline
For any proposition $P$
\newline
If $P(0)$, $\forall n \in \mathbb{N},P(n) \to P(n+1)$, then $\forall n \in \mathbb{N}, P(n)$
\newline
We say $P(0)$ is a base case.
\newline
Inductive Step: $P(n) \to P(n+1)$
\newline
Inductive Hypothesis: $P(n)$
\end{defn}

\begin{exmp}
Use induction to show that $\forall n\in \mathbb{N}, n< 2^n$.
\begin{proof}
Let $P(n)=:\text{" $ n< 2^n$"}$
\newline
We proceed by induction:
\newline
Base Case, $n=0$
\newline
Since $0<1$ we see $P(0)$ is true. 
\newline
Inductive Step:
\newline
Assume the Inductive Hypothesis (I.H.)
\newline
Goal:
\newline
We assume $P(n)=\text{"$n<2^n$"}$ is true
\newline
W.T.S. $P(n+1)=:\text{" $ (n+1)< 2^{(n+1)}$"}$
\newline
\begin{align}
(n+1)&=n+1\\
    &<2^n+1 \text{ By I.H.}\\
    &=2^n+2^0\\
    &\leq 2^n+2^n
    &\leq 2*(2^n)\\
    &=2^{n+1}
\end{align}
By induction we see that,
$\forall n \in \mathbb{N}, n<2^n$

\end{proof}

\end{exmp}
\begin{exmp}
Use induction to show that $n^3-n$ is divisible by 3 whenever $n$ is a positive integer. 
\begin{proof}
Let $P(n)=\text{"$n^3-n$ is divisible by 3"}$
\newline

\end{proof}
We proceed by induction
\newline
Base Case:
\newline
We see that $P(1)$ is true since 0 is divisible by 3.
\newline
Inductive Step: $P(n)\to P(n+1)$
\newline
Inductive Hypothesis: $P(n)$
\newline
Assume $n^3-n$ is divisible by 3.
\newline
W.T.S. $P(n+1)$
\newline
$(n+1)^3-(n+1)$ is divisible by 3.
\newline
\begin{align}
(n+1)^3-(n+1)=&(n+1)(n+1)(n+1)-(n+1)\\
=&n^3+3n^2+3n+1-(n+1)\\
=&n^3-n +3n^2+3n+1-(1)\\
=&n^3-n +3n^2+3n\\
=&n^3-n +3(n^2+n)
\end{align}
By Inductive Hypothesis $n^3-n$, and since $3(n^2+n)$ is divisible by three, their sum is divisible by 3 and we have show $P(n+1)$.
\end{exmp}
\newpage
